%%%%%%%%%%%%%%%%%%%%%%%%%%%%%%%%%%%%%%%%%%%%%%%%%%%%%%%%%%%%%%%%%%%%%%%%%%%%
% AGUJournalTemplate.tex: this template file is for articles formatted with LaTeX
%
% This file includes commands and instructions
% given in the order necessary to produce a final output that will
% satisfy AGU requirements, including customized APA reference formatting.
%
% You may copy this file and give it your
% article name, and enter your text.
%
% guidelines and troubleshooting are here: 

%% To submit your paper:
\documentclass[final]{agujournal2019}
\usepackage{url} %this package should fix any errors with URLs in refs.
\usepackage{lineno}
\usepackage[inline]{trackchanges} %for better track changes. finalnew option will compile document with changes incorporated.
\usepackage{soul}
% Preamble ------------------------------------------------
\usepackage{preamble}
% Macros --------------------------------------------------
\usepackage{macros}
% Symbols --------------------------------------------------
\usepackage{symbols}
% Images --------------------------------------------------
\graphicspath{{assets/figures/}}


\linenumbers
%%%%%%%
% As of 2018 we recommend use of the TrackChanges package to mark revisions.
% The trackchanges package adds five new LaTeX commands:
%
%  \note[editor]{The note}
%  \annote[editor]{Text to annotate}{The note}
%  \add[editor]{Text to add}
%  \remove[editor]{Text to remove}
%  \change[editor]{Text to remove}{Text to add}
%
% complete documentation is here: http://trackchanges.sourceforge.net/
%%%%%%%

\draftfalse

%% Enter journal name below.
%% Choose from this list of Journals:
%
% JGR: Atmospheres
% JGR: Biogeosciences
% JGR: Earth Surface
% JGR: Oceans
% JGR: Planets
% JGR: Solid Earth
% JGR: Space Physics
% Global Biogeochemical Cycles
% Geophysical Research Letters
% Paleoceanography and Paleoclimatology
% Radio Science
% Reviews of Geophysics
% Tectonics
% Space Weather
% Water Resources Research
% Geochemistry, Geophysics, Geosystems
% Journal of Advances in Modeling Earth Systems (JAMES)
% Earth's Future
% Earth and Space Science
% Geohealth
%
% ie, \journalname{Water Resources Research}

\journalname{Journal of Advances in Modeling Earth Systems (JAMES)}


\begin{document}

%%%%%%%%%%%%%%%%%%%%%%%%%%%%%%%%%%%%%%%%%%%%%%%
%  TITLE
%
% (A title should be specific, informative, and brief. Use
% abbreviations only if they are defined in the abstract. Titles that
% start with general keywords then specific terms are optimized in
% searches)
%
%%%%%%%%%%%%%%%%%%%%%%%%%%%%%%%%%%%%%%%%%%%%%%%

% Example: \title{This is a test title}

\title{Accounting for truncation error in reduced-gravity models: a subsurface-aware model for enhanced ocean surface reconstruction}

%%%%%%%%%%%%%%%%%%%%%%%%%%%%%%%%%%%%%%%%%%%%%%%
%
%  AUTHORS AND AFFILIATIONS
%
%%%%%%%%%%%%%%%%%%%%%%%%%%%%%%%%%%%%%%%%%%%%%%%

% Authors are individuals who have significantly contributed to the
% research and preparation of the article. Group authors are allowed, if
% each author in the group is separately identified in an appendix.)

% List authors by first name or initial followed by last name and
% separated by commas. Use \affil{} to number affiliations, and
% \thanks{} for author notes.
% Additional author notes should be indicated with \thanks{} (for
% example, for current addresses).

% Example: \authors{A. B. Author\affil{1}\thanks{Current address, Antartica}, B. C. Author\affil{2,3}, and D. E.
% Author\affil{3,4}\thanks{Also funded by Monsanto.}}

\authors{Gaétan Rigaut\affil{1}, Noé Lahaye\affil{1}, Etienne Mémin\affil{1}}


% \affiliation{1}{First Affiliation}
% \affiliation{2}{Second Affiliation}
% \affiliation{3}{Third Affiliation}
% \affiliation{4}{Fourth Affiliation}

\affiliation{1}{Centre INRIA de l'Université de Rennes, France}
%(repeat as many times as is necessary)


% Corresponding author mailing address and e-mail address:

% (include name and email addresses of the corresponding author.  More
% than one corresponding author is allowed in this LaTeX file and for
% publication; but only one corresponding author is allowed in our
% editorial system.)

% Example: \correspondingauthor{First and Last Name}{email@address.edu}

\correspondingauthor{Gaétan Rigaut}{gaetan.rigaut@inria.fr}



%%%%%%%%%%%%%%%%%%%%%%%%%%%%%%%%%%%%%%%%%%%%%%%
% KEY POINTS
%%%%%%%%%%%%%%%%%%%%%%%%%%%%%%%%%%%%%%%%%%%%%%%
%  List up to three key points (at least one is required)
%  Key Points summarize the main points and conclusions of the article
%  Each must be 140 characters or fewer with no special characters or punctuation and must be complete sentences

% Example:
% \begin{keypoints}
% \item	List up to three key points (at least one is required)
% \item	Key Points summarize the main points and conclusions of the article
% \item	Each must be 140 characters or fewer with no special characters or punctuation and must be complete sentences
% \end{keypoints}

\begin{keypoints}
    \item Enhancing reduced-gravity model reconstruction in the context of data assimilation by acknowledging the truncation error
    \item Introducing a physics-derived regularization term for the cost function to enforce multilayer QG dynamics
    \item Improving the robustness of the model by introducing a deformation radius correcting parameter in the control vector
\end{keypoints}

%%%%%%%%%%%%%%%%%%%%%%%%%%%%%%%%%%%%%%%%%%%%%%%
%
%  ABSTRACT and PLAIN LANGUAGE SUMMARY
%
% A good Abstract will begin with a short description of the problem
% being addressed, briefly describe the new data or analyses, then
% briefly states the main conclusion(s) and how they are supported and
% uncertainties.

% The Plain Language Summary should be written for a broad audience,
% including journalists and the science-interested public, that will not have 
% a background in your field.
%
% A Plain Language Summary is required in GRL, JGR: Planets, JGR: Biogeosciences,
% JGR: Oceans, G-Cubed, Reviews of Geophysics, and JAMES.
% see http://sharingscience.agu.org/creating-plain-language-summary/)
%
%%%%%%%%%%%%%%%%%%%%%%%%%%%%%%%%%%%%%%%%%%%%%%%

%% \begin{abstract} starts the second page

\begin{abstract}
      Ocean surface reconstruction through variational assimilation (VA) of partial, satellite-acquired observations is usually performed using reduced-order models. 
      The reduced-gravity (RG) framework, a truncation of the multilayer quasi-geostrophic (QG) equations, has shown promising reconstructing capabilities.
      In order to address model errors, a state-of-the-art reconstruction method includes an external forcing in the advection equations of the RG model and optimizes this forcing term during the VA procedure.
      In this paper, we show that the RG advection equation can be corrected to account specifically for the truncation error by inferring the material derivative of the subsurface stream function, which can be parameterized and optimized using VA.
      Since this method yields a proxy for the subsurface stream function, we can compute an approximation of the subsurface potential vorticity (PV) and an estimation of its conservation.
      We use this approximated conservation as a physics-based regularization term for the VA cost function, replacing the usual background regularization.
      We also introduce a single additional control parameter for deformation radius correction, allowing the model to mitigate incorrect stratification parameters.
      Method validation is performed using two different reference datasets, a three-layer QG double-gyre simulation and the eNATL60 simulation, both with observations retrieved by an OSSE along SWOT-like tracks.
      Results show a more than 80\% and 25\% reduction of surface stream function reconstruction error in each experiment (respectively) when compared to a RG model corrected by an external forcing.
      This study emphasizes the benefits of acknowledging simplification assumptions to design theory-informed error parameterizations to further refine reconstruction capabilities of reduced-order models.
\end{abstract}

\section*{Plain Language Summary}
Enter your Plain Language Summary here or delete this section.
Here are instructions on writing a Plain Language Summary: 
https://www.agu.org/Share-and-Advocate/Share/Community/Plain-language-summary


%%%%%%%%%%%%%%%%%%%%%%%%%%%%%%%%%%%%%%%%%%%%%%%
%
%  BODY TEXT
%
%%%%%%%%%%%%%%%%%%%%%%%%%%%%%%%%%%%%%%%%%%%%%%%

%%% Suggested section heads:
% \section{Introduction}
%
% The main text should start with an introduction. Except for short
% manuscripts (such as comments and replies), the text should be divided
% into sections, each with its own heading.

% Headings should be sentence fragments and do not begin with a
% lowercase letter or number. Examples of good headings are:

% \section{Materials and Methods}
% Here is text on Materials and Methods.
%
% \subsection{A descriptive heading about methods}
% More about Methods.
%
% \section{Data} (Or section title might be a descriptive heading about data)
%
% \section{Results} (Or section title might be a descriptive heading about the
% results)
%
% \section{Conclusions}



\section{Introduction}

In physical oceanography, the majority of available data is acquired by satellites that cannot measure the interior ocean state. However, the constellation of space-borne altimetric sensors is not dense enough to provide a full coverage of the globe. In order to provide estimations of the ocean state from partial observations, reconstructions method have been investigated since the first altimetric mission was launched in the late 1970s. Among them, the optimal interpolation method~\cite{rienecker1991,traon1998}, a purely statistical method that solely uses data was applied to reconstruct surface state. This method, as of today, is still at the basis of the DUACS product which remains a widely used real-time gridded ocean product.

In order to combine both simulated and observed data while accounting for their specific sources of error~\cite{lahoz2010}, data assimilation (DA) techniques started to be applied to the ocean system. Integrating observations into high-resolution numerical simulations is generally computationally intensive as it often requires multiple runs of a numerical model. In order to efficiently perform DA, observations are often assimilated into so-called reduced-order models (models that rely on simplified frameworks to emulate ocean system dynamics), although not exclusively~\cite{mogensen2012,cooper1996}.

Around the midlatitudes, a simplified framework of choice is the quasi-geostrophic (QG) theory. This framework, developed in the 1950s~\cite{charney1950,charney1953}, describes the behavior of flows around the geostrophic balance through the material conservation of potential vorticity (PV) by the geostrophic flow. Such theory has paved the way for many simplified numerical models, whether atmospheric or oceanic~\cite{phillips1951,held1995}. Although it describes 3D dynamics, it can be discretized by stacking layers of fluid with constant vertical properties, yielding the multilayer QG framework~\cite{vallis2006}. 

Studies in the late 1980s have used QG systems to perform data assimilation. Some would consist in updating the surface state with observations and waiting for the model's dynamics to adjust to its observation-imposed circulation~\cite{haines1991}. However, in order to produce satisfying reconstructions, these methods would require extensive spin-up times and thus significant amount of data. Other methods focused on using statistical information on the flow dynamics to project surface observations onto empirical vertical modes~\cite{mey1987}. Despite being promising, they would require many subsurface observations to infer relevant statistics thus preventing from extrapolating such techniques to the wide ocean.

Among DA methods, variational assimilation (VA) has been successfully employed to reducing observations to numerical simulation mismatches~\cite{gunson1996,cong1998,zou1995,vogeler1995,vogeler1999}. VA consists in performing successive iterations of a numerical model, from which reconstruction error is assessed. The backpropagation of the error using the adjoint model of the numerical code~\cite{errico1997} then allows refining the control parameters (initial condition for example) from an \textit{a priori} estimation to reduce such error. Despite the (significant) computation cost of determining and storing the adjoint model, VA is still widely used in operational assimilation softwares~\cite{martin2025}.

Performing VA using a multilayer QG simulation has been widely investigated in the late 1980s~\cite{cong1998,vogeler1995,vogeler1999}. But it has been pointed out~\cite{vogeler1995,vogeler1999,hurlburt1986,morrow1995} that, in the absence of direct subsurface observations, an accurate first guess of the initial condition was required to further refine surface dynamics reconstruction.

With the recent launch of the Surface Water and Ocean Topography (SWOT) mission~\cite{fu2024} has come the perspective of sea surface reconstruction with unprecedented resolution~\cite{benkiran2026}. This satellite provides data over wide swaths, but its 21-day repeat cycle results in significant gaps in the collected data, emphasizing the need for interpolation techniques. Therefore, both the lack of detailed subsurface direct observations and the motivation of only reconstructing surface flow have motivated the use of reduced-gravity (RG) QG models in data assimilation. This framework is directly derived from the multilayer one by considering deep ocean flow to be at rest. This additional simplification offers an evolution equation for the surface flow without requiring any information on the subsurface state. 

Using a RG model,~\citeA{ubelmann2015} presented a so-called dynamical interpolation method which refines optimal interpolation reconstruction by benefitting from the ability of the model to render non-linear dynamics~\cite{ubelmann2016}. In order to account for any source of error in the RG model, the VarDyn framework~\cite{leguillou2023,leguillou2025} introduces an external forcing into the RG advection equation. Such forcing is estimated through VA to approach the rich and complex dynamics of the ocean surface. However, in both methods remains the approximation that the flow is at rest below the ocean surface.

In the present study, we rewrite the top layer equation of the multilayer QG framework to derive a truncation-aware RG equation. This formulation is based on the introduction of an additional forcing composed of the sublayer stream function's material derivative. VA is then performed to find the sublayer field that allows for the most accurate reconstruction of the surface flow. As our method provides a direct estimation for the sublayer stream function, we are also able to replace the usual background term of the cost function by a physics-based regularization enforcing PV conservation in the sublayer. Additionally, we introduce a single-parameter deformation radius correction to help the model handling incorrect stratification parameters.
In order to validate our method, we perform two experiments. We first use a three-layer QG double-gyre simulation as reference to monitor truncation-only error mitigation and discuss the relevancy of the parameterized field as a proxy for the sublayer stream function. Then, we use a primitive equation simulation to assess the model's reconstruction capabilities of a flow with richer dynamics.



\section{Reduced-gravity QG model with parameterized sublayer flow}
In this section, we present a modified reduced-gravity QG model that aims at capturing the surface currents while taking into account the sub-surface dynamics. 
We begin with a reminder of the standard multilayer Quasi-Geostrophic theory and the truncated one-layer reduced-gravity model. 

\subsection{Multilayer Quasi-Geostrophic theory}
QG dynamics relies on the material conservation of PV by geostrophic flow. Both the PV and the geostrophic flow can be expressed using the stream function which (through the inversion of the elliptic operator mapping stream function to PV) allows recovering the geostrophic flow from the PV itself~\cite{hoskins1985}.  Multilayer QG, the vertical discretization of the continuously stratified QG equation, is obtained by stacking layers of fluid with constant vertical properties~\cite{pedlosky1987,vallis2006} inside which PV is conserved. In the following we will recall these equations in a $N$-layer fluid. Using $\vect{\cdot}$ to represent layered fields ($\vect{\chi}= \transpose{\begin{bmatrix}\chi_1, \dots, \chi_N\end{bmatrix}}$), we denote
\begin{itemize}
    \item $\vect{\rho}$ the reference layer densities,
    \item $\vect{H}$ the reference layer thickness,
    \item $\vect{g}$ the reduced gravities ($g_i = g\left(\rho_{i}-\rho_{i-1}\right)/\rho_1$, with the convention $g_1 = g$),
    \item $\vect{\sf}$ the stream function vector,
    \item $\vect{\pv}$ the PV vector.
\end{itemize}

Introducing the horizontal gradient $\grad = \transpose{\begin{bmatrix} \pdvx & \pdvy \end{bmatrix}}$ along with the perpendicular horizontal gradient $\gradperp = \transpose{\begin{bmatrix} -\pdvy & \pdvx \end{bmatrix}}$, the free-surface multilayer QG equations (in the absence of external forcing) in the layer $i$ read
\begin{subequations} \label{eq:mlqg}
    \begin{align}
        \pv_i = \lap \sf_i - f_0^2 {\left(\mat{A} \vect{\sf}\right)}_i, \label{def:pv-ml}\\
        \pdvt{\pv_i} + \gradperp \sf_i \cdot \grad \pv_i = 0\label{def:pv-adv}.
    \end{align}
\end{subequations}
Here, the stretching matrix $\mat{A}$, is given by
\begin{equation}
    \mat{A} = \begin{bmatrix}
            \frac{1}{H_1 g_1} + \frac{1}{H_1 g_2} & -\frac{1}{H_1 g_2} & \cdot & \cdot \\
            -\frac{1}{H_2 g_2} & \frac{1}{H_2 g_2} + \frac{1}{H_2 g_3} & -\frac{1}{H_2 g_3} & \cdot \\
            \vdots & \ddots & \ddots & \vdots \\ 
            \cdot & \cdot & -\frac{1}{H_N g_N} & \frac{1}{H_N g_N} \\ 
        \end{bmatrix}.
\end{equation}

One sees that the different layers are coupled with each others via the vortex stretching term in the equation for the PV. 
In particular, the upper-layer PV is coupled with the second layer, as follows from its definition:
\begin{equation}\label{def:pv-top-ml}
    \pv_1 = \lap{\sf_1} - \frac{f_0^2}{H_1 g_1} \sf_1 - \frac{f_0^2}{H_1 g_2} \left(\sf_1 - \sf_2\right).
\end{equation}

From the above set of equations \eqref{eq:mlqg}, one can derive an approximate single-layer (so-called reduced-gravity) model by truncating the number of layers to 1 and ignoring the coupling with $\sf_2$ in the PV equation.
The latter simplification is strictly equivalent to assuming $\sf_2=0$, \textit{i.e.} a lower layer at rest. 
The corresponding top-layer PV, denoted $\pv^{RG}$ throughout this paper, thus becomes
\begin{equation}\label{def:pv-rg}
    \pv^{RG} = \lap{\sf_1} - f_0^2 \left(\frac{1}{H_1 g_1} + \frac{1}{H_1 g_2}\right) \sf_1,
\end{equation}
and the reduced-gravity QG equation, which describes the material conservation of this quantity transported by the flow $\sf_1$, is given by
\begin{equation} \label{eq:qg-rg}
    \pdvt{\pv^{RG}} +  \gradperp{\sf_1}\cdot \grad{\pv^{RG}} = 0.
\end{equation}
The prefactor in the vortex-stretching term in~\cref{def:pv-rg} defines the (inverse square of the) equivalent Rossby deformation radius (wherein the first term is usually negligible since $g_1\gg g_2$).
In practice, instead of determining typical layer depths and reduced gravity constants ($H_1, g_1$, etc.), this prefactor is usually replaced using the deformation radius associated with the first baroclinic mode, that is the second eigenvalue of the stretching matrix $\mat{A}$.

\subsection{Reduced-gravity QG model for the assimilation of altimeter data}
Direct observations through satellite sensing are only able to measure sea surface height (SSH), which can be matched to the top layer stream function $\sf_1$ in the QG framework. 
Remote sensing alone can not provide direct measurement on the sublayer flow $\sf_2$. 
As a result, the reduced-gravity model as expressed in~\cref{eq:qg-rg} was used in the first time to perform dynamical interpolation~\cite{ubelmann2015,ubelmann2016}. 
Benefiting from the ability of the QG model to generate non-linear dynamics, such method has been able to outperform statistics-based interpolation. 
However, due to the truncated QG equations, significant errors can emerge when performing optimal interpolation using solely the RG model, as pointed out in~\citeA{manucharyan2021}. 
In order to address this limitation, the VarDyn framework~\cite{leguillou2023,leguillou2025} takes into account the model error by introducing an external forcing, $\epsilon$, into the advection equation,
\begin{equation}
    \pdvt{\pv^{RG}} + J\left(\sf_1, \pv^{RG}\right) = \epsilon.
\end{equation}

Here, we propose a model for the error by assuming that this error term describes the impact of the lower-layer dynamics on the surface layer, following the derivation of the reduced-gravity QG equations above. 
Indeed, given the QG PV definition, the top layer PV can be written
\begin{equation}
    \pv_1 = \pv^{RG} + \frac{f_0^2}{H_1 g_2} \sf_2.
\end{equation}
Using this decomposition into the top layer PV advection equation~\eqref{def:pv-adv} yields
\begin{subequations}\label{def:qg-rg-adv-corr}
    \begin{align}
        \pdvt{\pv^{RG}} + J\left(\sf_1, \pv^{RG}\right) & = - \frac{f_0^2}{H_1 g_2} \left[\pdvt{\sf_2} + \gradperp \sf_1 \cdot \grad \sf_2 \right] \\
        & = - \frac{f_0^2}{H_1 g_2} \Dt{\sf_2}.
    \end{align}
    
\end{subequations}
From this expression, we note that the usual RG model advection matches the multilayer formulation when $\Dt*{\sf_2}=0$. 
Addressing the truncation error of the RG model resumes to emulating the right-hand side of~\cref{def:qg-rg-adv-corr}. 
Since $\sf_1$ is the prognostic variable of the model, computing $\Dt*{\sf_2}$ boils down to estimating $\sf_2$. Denoting $\tilde{\sf}_2$ such estimation, we define the corrected RG advection by
\begin{equation}
    \pdvt{\pv^{RG}} + J\left(\sf_1, \pv^{RG}\right) = - \frac{f_0^2}{H_1 g_2} \Dt{\tilde{\sf}_2}.
\end{equation}

In the following, this model will be referred to as RG-SI (Reduced Gravity -Sublayer Inferred). Additionally to explicitly addressing the RG-related truncation error, this method also provides a proxy for the sublayer stream function, through $\tilde{\sf}_2$. Our specific parameterization of $\tilde{\sf}_2$ will be detailed in~\cref{sec:forward-model}. 

\subsection{Further extensions and remarks on the Reduced-Gravity Sublayer-Inferred QG model}
\subsubsection{Approximate lower-layer PV conservation} \label{sec:PV2reg}
The error model introduced above for the upper layer PV equation relies on an estimate of the sublayer stream function $\sf_2$.
By construction, this variable cannot be determined as precisely as the surface stream function, and can lead to ill-constrained inversion. 
One strategy to mitigate this issue (in addition to a reduced-order formulation introduced later in~\cref{sec:ROpsi2}) is to rely on an approximate conservation law, namely conservation of the lower layer PV.
Following~\cref{def:pv-adv}, PV advection in the layer 2 reads
\begin{equation}
    \pdvt*{\pv_2} + \gradperp \sf_2 \cdot \grad \pv_2 = 0, 
    \quad q_2 = \lap{\tilde{\sf}_2} - \frac{f_0^2}{H_2 g_2}\left(\tilde{\sf}_2 - \sf_1\right) + \frac{f_0^2}{H_2 g_3}\left(\tilde{\sf}_3 - \sf_3\right)
\end{equation}
Even if one knows the stream function in the first two layers ($\sf_1$ and $\sf_2$), the system of equation is not closed as the second layer is coupled to $\sf_3$ -- which remains unknown -- via the vortex stretching terms in the PV definition.
We thus resort to a truncated sublayer PV definition (setting $\sf_3 = 0$) to obtain a proxy for $\pv_2$ and its evolution:
\begin{equation}
    \tilde{\pv}_2 = \lap{\tilde{\sf}_2} - \frac{f_0^2}{H_2 g_2}\left(\tilde{\sf}_2 - \sf_1\right).
\end{equation}
Since in our approach, $\sf_2$ will remain approximated, this equation will only be used in a weak sense to constrain the optimization process based on physical consideration (see below).

\subsubsection{Deformation radius correction}\label{sec:deformation-radius-correction}
Other than the truncation error, incorrect model stratification parameters (the $H_i$ or the $g_i$) could result in mismatch in the surface flow dynamics. Since a key characteristic of the flow dynamics is first baroclinic deformation radius, in order to increase the robustness of our two-layer-like model, we wish to allow for adjustments of such radius. Given the diagonalization of the two-layer stretching matrix,

\begin{equation}
    \mat{A}= \mat{C}_{\mathrm{M2L}}\begin{bmatrix}\lambda_{\rightleftarrows} & 0 \\ 0 & \lambda_{\rightrightarrows}\end{bmatrix}\mat{C}_{\mathrm{L2M}},
\end{equation}

where $\lambda_{\rightleftarrows}$ is the baroclinic-related eigenvalue and $\lambda_{\rightrightarrows}$ the barotropic one, and $\mat{C}_{\mathrm{M2L}}$ (resp., $\mat{C}_{\mathrm{L2M}}$) the modes to layers (resp., layers to modes) passage matrix. The baroclinic radius correction is performed by introducing the parameter $\alpha$, and the modified stretching matrix,

\begin{subequations}
    \begin{align}
        \tilde{\mat{A}}_{\alpha} & = \mat{C}_{\mathrm{M2L}}\begin{bmatrix}\left(1+\alpha\right)\lambda_{\rightleftarrows} & 0 \\ 0 & \lambda_{\rightrightarrows}\end{bmatrix}\mat{C}_{\mathrm{L2M}}, \\
        & = \mat{A} + \alpha\lambda_{\rightleftarrows}\mat{C}_{\mathrm{M2L}}\begin{bmatrix}1 & 0 \\ 0 & 0\end{bmatrix}\mat{C}_{\mathrm{L2M}},
    \end{align}
\end{subequations}

which coincides with the default stretching matrix when $\alpha=0$. The barotropic deformation radius then remains unchanged as it is linked to the total depth which can be quite easily approximated whereas the baroclinic deformation radius, $L_{\rightleftarrows}$ becomes

\begin{equation}
    L_{\rightleftarrows} = \frac{1}{f_0 \sqrt{\left(1+\alpha\right)\lambda_{\rightleftarrows}}}.
\end{equation}

Our model implementation replaces $\mat{A}$ by $\tilde{\mat{A}}_{\alpha}$ in the usual two-layer QG equations and the parameter $\alpha$ is added in the control vector, granting one additional degree of freedom for the variational assimilation.

\section{Variational assimilation of altimeter data with the RG-SI model}
This section describes how the model presented above is implemented in a variational assimilation framework for inverting surface altimeter data, starting with a brief presentation of the cost function used, and followed by a description of the various modeling strategies applied to the RG-SI model.

\subsection{General variational formulation}\label{sec:va}
Provided a set of (partial) observations of the surface stream function $\left( \Psi_i \right)$ retrieved at times $\left( t_i \right)$, $i \in \{1, \dots, N\}$, we emulate the surface state using a model $\mathcal{M}_{\theta}$ (RG-SI for example), which depends on the control vector $\theta$. The latter may include any parameter the propagator dynamics depends on (initial and boundary conditions, model parameters\dots). Estimated states are then processed through the observation operator $\mathcal{H}$ in order to be compared to the observations. Finally, variational assimilation reduces to finding $\theta^\star$ such that it minimizes the observation cost function, hence

\begin{equation}
    \theta^\star = \arg \min_{\theta} \sum_{i=1}^{N}{\lVert \Psi_i - \mathcal{H}_i \circ \mathcal{M}_{\theta}\left( t_i\right)\rVert}^2_{\mat{R}^{-1}} = \arg \min_{\theta} \loss{obs}\left(\theta\right),
\end{equation}

with the observation error covariance matrix $\mat{R}$ (in the experiments from~\cref{sec:method-validation}, this matrix is taken as the identity, up to a multiplicative factor for normalization purpose).

However, there are no guarantees that enough data is available for the minimization problem to be well constrained. Therefore, provided a prior estimate of the control vector, $\theta^b$ (called the background), and its error covariance matrix, $\mat{B}$, the cost function is usually regularized by penalizing the deviation of $\theta$ from $\theta^b$. Introducing the so-called background regularization cost function
\begin{equation}
    \loss{reg}\left(\theta\right) = {\lVert \theta - \theta^b \rVert}^2_{\mat{B}^{-1}},
\end{equation}
the cost function rewrites
\begin{equation}
    \loss\left(\theta\right) = \loss{obs}\left(\theta\right) + \loss{reg}\left(\theta\right).
\end{equation}

Retrieving a relevant estimation for $\theta^b$ can already be a challenge of itself and should not be overlooked. Indeed, when performing data assimilation using multilayer QG models, it is often difficult to obtain a precise \textit{a priori} estimation for the subsurface flow that can be used as a background term. To overcome this limitation, one could penalize deviation from 0 or from the initial surface state, $\sf_1^0$, but doing so would strongly call back the initial condition toward either a surface intensified or a barotropic flow. Since this background term is solely there to ensure the well-posedness of the minimization problem, the background regularization can be replaced by any cost function that would enforce an \textit{a priori} knowledge on the dynamics. For instance, penalizing the stream function's second derivatives from deviating from 0 has been proven to be effective in regularizing the assimilation problem as it enforces the smoothness of the fields~\cite{vogeler1995,cong1995}.

However, this does not replace the need for an actual initial condition to start assimilating from. To overcome this difficulty,~\citeA{cong1998} decided to use a preparation step, prior to performing variational assimilation, in order to find a suitable linear combination of $\sf_1^0$ and $0$ to instantiate the model with. Still, not only does an appropriate initial condition help the assimilation algorithm converges towards a minimum~\cite{morrow1995}, it also helps the model in emulating relevant dynamics. Indeed, as shown by~\citeA{hurlburt1986}, instantiating a two-layer primitive equation model using a RG-like condition (hence inert bottom layer at instantiation) yields little to no better results than using the equivalent one-layer reduced-gravity primitive equation model. Hence, in the absence of model error correction, increasing the model's complexity does not suffice to mitigate inaccurate initial conditions.

\subsection{Specific formulation of the RG-SI model parameters for data assimilation}\label{sec:forward-model}
Since the focus of this paper is to address the unknown sublayer dynamics, we restrict ourselves to estimating parameters that connect only to this variable in the optimization. 
In particular, initial conditions, which are typically inferred in data assimilation, along with boundary conditions, are assumed known and their inversion is not addressed. 
Thus, we only need to address how the terms related to the sublayer dynamics are to be implemented in an efficient way.

\subsubsection{Reduced-order formulation for $\tilde{\sf}_2$}\label{sec:ROpsi2}
In order to keep a reasonable number of parameters for the expression of $\tilde{\sf}_2$, it is emulated using a reduced-order space. 
We express $\tilde{\sf}_2$ over a quite simple space -- namely a partition of unity of space-time Gaussian envelopes --, so that RG-SI performances could not be attributed to the structure of the parameterization rather than the actual error model.
We denote $N_T$ the number of time-related envelopes and $N_{XY}$ the number of space-related envelopes at a given time $t_i$. 
Given $i \in \{1,\dots,N_T\}$ and $j \in \{1,\dots,N_{XY}\}$, we define
\begin{subequations}
    \begin{align}
        e^T_{i}\left(t\right) & = \exp\left[-{\left(\frac{t-t_i}{\sigma_T}\right)}^2\right], \\
        e^{XY}_{j}\left(x,y\right) & = \exp\left[-{\left(\frac{x-x_j}{\sigma_X}\right)}^2\right]\exp\left[-{\left(\frac{y-y_j}{\sigma_Y}\right)}^2\right],
    \end{align}
\end{subequations}
where $t_i, x_j, y_j$ are the locations of the envelope centers and $\sigma_T, \sigma_X, \sigma_Y$ their standard deviations. Envelopes centers are chosen so that they are evenly spaced (whether in time or space). Details on the determination of the standard deviations can be found in~\cref{appendix:rg-si-parameterization}.
We denote $E^T = \sum_{i=1}^{N_T} e^T_i$ and $E^{XY} = \sum_{i=1}^{N_{XY}} e^{XY}_j$. Using these notations, the overall decomposition of $\tilde{\sf}_2$ at time $t$ and location $\left(x,y\right)$ reads
\begin{equation}
    \tilde{\sf}_2\left(t,x,y\right) = \frac{\sum_{i=1}^{N_T}\sum_{j=1}^{N_{XY}}  \theta_{i,j} e^T_{i}\left(t\right)e^{XY}_{j}\left(x,y\right)}{E^T\left(t\right) E^{XY}\left(x,y\right)},
\end{equation}
where the ${\left(\theta_{i,j}\right)} \in \mathbb{R}^{N_T \times N_{XY}}$ are coefficients inferred through cost function minimization. 

This decomposition allows representing a smooth field while enabling the emergence of some rather small scale patterns. 
Additionally, we can easily restore the dynamics of a flow under the RG assumption by setting all coefficients to 0. 
In the forthcoming, this will be the initial condition considered when instantiating the variational assimilation procedure.

\subsubsection{Physics-based regularization}
Efficiency of the optimization is further improved by using a physic-based regularization based on an approximate form of the conservation of PV, as introduced in~\cref{sec:PV2reg}.
The corresponding regularization term reads
\begin{equation}
    \loss{reg}\left(\sf_1, \tilde{\sf}_2\right) = {\lVert \pdvt*{\tilde{\pv}_2} + \gradperp \tilde{\sf}_2 \cdot \grad \tilde{\pv}_2 \rVert}^2.
\end{equation}
The overall cost function is thus defined as
\begin{equation}\label{def:RG-SI-reg}
    \loss \left(\tilde{\sf}_2\right) = \loss{obs}\left(\tilde{\sf}_2\right) + \gamma \loss{reg}\left(\tilde{\sf}_2\right),
\end{equation}
where $\gamma$ is an hyperparameter driving the relative importance of the regularization over the observation loss. 

Replacing the usual background state related regularization by this physics-based formulation is inspired by physics-informed neural networks (PINN) methodology~\cite{raissi2019} and PINN-derived data assimilation approaches have proven the relevancy of regularizing using physical constraint~\cite{ugur2025,oh2025}.

As explained in~\cref{sec:va}, switching to a physics-based regularization removes the need for an a priori guess of the parameter. Instead of recalling $\tilde{\sf}_2$ toward 0 (since it would be the natural choice in the absence of subsurface flow data), it allows the model to stir away from the RG framework. Rather than constraining the value of the sublayer stream function itself (as a background regularization on its initial condition would do), such regularization enforces two-layer-like motion. Only does the model need to satisfy a pseudo-advection of PV in the bottom layer, but as for its value it remains unconstrained.

It is worth emphasizing that introducing this regularization does not strictly enforce the conservation of PV in the sublayer. It guides the model towards finding an expression for $\tilde{\sf}_2$ that approaches the conservation $\tilde{\pv}_2$. However, as there are no guarantees that $\tilde{\pv}_2$ should be strictly conserved ($\tilde{\sf}_2$ is not exact and $\tilde{\pv}_2$ is computed through the truncation of the multilayer equations), this formulation allows for some flexibility on said conservation. Indeed, if improving surface dynamics requires drifting away from the advection of the sublayer (truncated) PV (through the introduction of non-QG patterns in the reconstructed field for example), it can do so. A strictly enforced two-layer QG model would not be able to do so without the introduction of an external forcing term in the sublayer advection equation.

\subsection{Numerical implementation}
The QG models described above, both one layer and multilayer, were implemented in Pytorch following the code of~\citeA{thiry2024} (extended to withstand non-zero boundary conditions). The corresponding Python code is available at \url{https://github.com/grigaut/QG/}.
All the numerical experiments presented below, except for the NEMO-based eNATL60 numerical simulation, were run using this implementation. 
We leveraged the library's built-in automatic differentiation~\cite{paszke2017} to both retrieve cost function gradient and perform Adam algorithm for finding the minimum of the cost function~\cite{kingma2017}.

Numerical experiments were conducted using the Grid'5000 testbed~\cite{balouek2013}.

\section{Method validation}\label{sec:method-validation}

We have considered two observing system simulation experiment (OSSE) validation experiments. In the first one, we have set a double-gyre experiment using a three-layer QG model (QG3L). The reconstruction error from the RG model (dynamical interpolation situation) is mostly explained by the truncation of the QG equations. Within this framework, we can assess the ability of the RG-SI model in reconstructing top layer dynamics and compare its performances with a RG model with no error correction (RG0) and a RG model whose error is prescribed though an external forcing, RG-F (RG-Forced). Two variants of this setup have been explored. One where model parameters (reference layer depth, reduced gravity values\dots) are the same in the reference QG model and in the RG counterparts. Not only does this configuration confine the error sources to truncation error, it also allows for a fair comparison of $\tilde{\sf}_2$ and $\sf_2$. In the other variant, model parameters are modified to assess the robustness of each of the models when dealing with incorrect baroclinic deformation radius. The second validation experiment uses reference surface fields from the high-resolution, realistic simulation eNATL60 simulation~\cite{brodeau2020}, which is based on NEMO~\cite{madec2024} an includes realistic forcings and bathymetry. This experiment investigates the capability of our model in depicting a flow that present ageostrophic, complex and rich dynamical features.

\subsection{Configuration of OSSEs}

As this work focuses on introducing a new parameterization to reduce model error, we have decided to only consider such parameterization in the control vector. Indeed, since all tested frameworks (RG0, RG-F, RG-SI) share identical initial and boundary conditions, differences in performances only result from model error and its correction. Ideal initial and boundary conditions are therefore retrieved from the surface flow of the reference simulation. Furthermore, the estimation of initial and boundary conditions through variational assimilation has already been widely investigated~\cite{cong1998,zou1995,gunson1996,seiler1993,ayoub2006}, suggesting that our methodology could quite easily be adapted to a more realistic data assimilation configuration.

In both experiments, observations are drawn from the surface using an OSSE that simulates SWOT-like sampling tracks. Satellite track example can be found in~\cref{fig:satellite-track}, along with the cumulative passages of the tracks over a 20-days period.

\begin{figure}[H]
    \begin{center}
        \includesvg[width=0.95\textwidth]{assets/images/satellite_track.svg}
    \end{center}
    \caption{Example of simulated satellite track over the surface (Surface fields extracted from a QG3L double-gyre simulation). Last panel shows how many times has a pixel been observed over the 20-days period.}\label{fig:satellite-track}
\end{figure}

Assimilation is conducted on 20 days-long periods. Over these periods, satellite passages are simulated and, when a track spans over the considered domain, observation loss is evaluated using the collected data. Regarding the RG-SI model, regularization cost function is computed at every time step. Details on the implementation of the RG-F forcing and regularization are detailed in~\cref{appendix:rg-f-implementation}.

During the variational assimilation process, all parameterized fields are instantiated at 0. Therefore, RG-SI's initial observation loss is the exact same as the RG0 one since the bottom layer appears to be at rest. Optimization was stopped after 200 iterations as longer optimization did not yield any significant improvement in the cost function for each of the models.

After the assimilation process is completed, model performances are evaluated using the normalized Root Mean Squared Error (nRMSE). Given a stream function field $\sf$ and its reference value $\sf_{\mathrm{ref}}$, the corresponding nRMSE is

\begin{equation}
    \mathrm{nRMSE}\left(\sf, \sf_{\mathrm{ref}}\right) = \sqrt{\frac{{\lVert \sf - \sf_{\mathrm{ref}} \rVert}^2}{{\lVert \sf_{\mathrm{ref}} \rVert}^2}}.
\end{equation}

Therefore, perfect field reconstruction would result in a nRMSE score of 0 whereas a nRMSE score above 1 indicates that no useful information on $\sf_{\mathrm{ref}}$ can be retrieved from $\sf$. In addition to nRMSE over stream function, we evaluate the nRMSE over the surface velocity $\symbf{u} = \transpose{\begin{bmatrix} u & v\end{bmatrix}}$, which given reference surface velocity $\symbf{u}_{\mathrm{ref}} = \transpose{\begin{bmatrix} u_{\mathrm{ref}} & v_{\mathrm{ref}} \end{bmatrix}}$ reads

\begin{equation}
    \mathrm{nRMSE}_{\mathrm{vel}}\left(\symbf{u}, \symbf{u}_{\mathrm{ref}}\right) = \sqrt{\frac{{\lVert u - u_{\mathrm{ref}} \rVert}^2 + {\lVert v - v_{\mathrm{ref}} \rVert}^2}{{\lVert u_{\mathrm{ref}} \rVert}^2 + {\lVert v_{\mathrm{ref}} \rVert}^2}}.
\end{equation}


\subsection{3-layer QG double gyre}
\subsubsection{Experimental set-up}

All models were configured with the same time and space resolutions. Regarding stratification parameters, RG0 has been implemented as a one-layer QG model with $g = g_1 g_2/(g_1+g_2)$ so that it matches the parameters of a truncation of the multilayer equations (see~\cref{def:pv-rg}). RG-SI and RG-F have been instantiated so that they share the same deformation radius, namely $\qty{29}{\kilo\meter}$, which corresponds to the baroclinic deformation radius of the two-layer QG model matching the parameters of the reference three-layer QG on the top two-layers. Such deformation radius value is also similar to the one described in~\citeA{leguillou2023}. All model parameters are detailed in~\cref{tab:qg3l-model-parameters}.

\begin{table}[H]
    \centering
    \begin{tabular}{lcccc}
        \hline
        $\dt$ & \multicolumn{4}{c}{$\qty{7200}{\second}$} \\
        $\dx$ & \multicolumn{4}{c}{$\qty{10}{\kilo\meter}$} \\
        $\dy$ & \multicolumn{4}{c}{$\qty{10}{\kilo\meter}$} \\
        $f_0$ & \multicolumn{4}{c}{$\qty{9.375e-5}{\per\second}$} \\
        $\beta$ & \multicolumn{4}{c}{$\qty{1.754e-11}{\per\meter\per\second}$} \\
        \midrule
         & QG3L & RG-SI & RG-F & RG0 \\
        \midrule
        $H_1$ & $\qty{400}{\meter}$ & $\qty{400}{\meter}$ & $\qty{293}{\meter}$ & $\qty{400}{\meter}$ \\
        $H_2$ & $\qty{1100}{\meter}$ & $\qty{1100}{\meter}$ &  &  \\
        $H_3$ & $\qty{2600}{\meter}$ &  &  &  \\
        $g_1$ & $\qty{9.81}{\meter\per\second^2}$ & $\qty{9.81}{\meter\per\second^2}$ & $\qty{0.025}{\meter\per\second^2}$ & $\qty{0.0249}{\meter\per\second^2}$ \\
        $g_2$ & $\qty{0.025}{\meter\per\second^2}$ & $\qty{0.025}{\meter\per\second^2}$ &  &  \\
        $g_3$ & $\qty{0.0125}{\meter\per\second^2}$ &  &  &  \\
        \midrule
        $R_d$ & $\qty{41}{\kilo\meter}$ & $\qty{29}{\kilo\meter}$ & $\qty{29}{\kilo\meter}$ & $\qty{34}{\kilo\meter}$ \\
        \hline
    \end{tabular}
    \caption{Model parameters. $R_d$ denotes the first baroclinic deformation radius.}\label{tab:qg3l-model-parameters}
\end{table}

Since the reference simulation is forced by wind to display a double-gyre, we also include the wind as an external forcing within the advection equations for all models. Denoting $\tau = \transpose{\begin{bmatrix} \tau_x & \tau_y \end{bmatrix}}$ the wind stress tensor, the advection equation rewrites

\begin{equation}
    \pdvt{q^{RG}} + \gradperp \sf_1 \cdot \grad {\pv}^{RG} = F + \gradperp \cdot \tau
\end{equation}

where $F$ represent the additional forcing that depends on the framework,

\begin{subequations}
    \begin{align}
        F_{\mathrm{RG}} & = 0, \\
        F_{\mathrm{RG-F}} & = \epsilon, \\
        F_{\mathrm{RG-SI}} & = \frac{-f_0^2}{H_1 g_2}\Dt{\tilde{\sf}_2}.
    \end{align}
\end{subequations}

As the reference was established over a $\qty{2560}{\kilo\meter}$ by $\qty{5120}{\kilo\meter}$ domain, we have decided to focus on subdomains of the said area. We selected two $\qty{640}{\kilo\meter}$ x $\qty{1280}{\kilo\meter}$ zones, Z1, Z2, presenting different dynamical features.

\begin{figure}[H]
    \centering
    \begin{center}
        \includesvg[width=0.7\textwidth]{assets/images/zones.svg}
    \end{center}
    \caption{Detail of the 4 zones of the three-layer QG experiment.}
\end{figure}

Z1 is at the core of the eastward jet and presents a surface intensified dynamics whereas on Z2 the barotropic component of the flow is stronger. In both zones, the OSSE follows the same tracks, observing 13243 gridpoints. Number of control parameters for the different models are detailed in the following table

\begin{table}[H]
    \centering
    \begin{tabular}{lcccc}
          & RG-SI & RG-SI ($\alpha = 0$) & RG-F & RG0 \\
        \hline
        Control vector size &  6322 & 6321 & 8512 & 0 \\
        \hline
    \end{tabular}
    \caption{Comparison of the control vector size for the different models in the QG3L experiment.}\label{tab:qg3l-control-vector-size}
\end{table}

Both RG-F and RG-SI regularizations are scaled through an hyperparameter ($\gamma$). The value of $\gamma$ has been taken constant over all areas and assimilation periods. Regarding RG-SI, its value has been fitted so that regularization cost function on first assimilation step was around 1\%-1‰ of the observation loss in all areas and periods. Therefore, cost function is dominated by the observation error at first, which is then reduced in priority. Regarding RG-F, the value $\gamma$ has been empirically chosen so that the model would perform as best as possible while not displaying non-physical small-scale artefacts in the reconstructed fields (which it tends to do in the absence of regularization).

providing fields that would present  small-scale intensified (which it tends to do in the absence of regularization).

\subsubsection{Surface stream function reconstruction evaluation}

Performances were assessed over 12 20-days periods, each separated by 20 days. For both Z1 and Z2, error scores were then averaged and the empirical standard deviation of the error computed and displayed as a dispersion around the average error as shown in~\cref{fig:qg3L-nrmses}. Velocity nRMSEs were computed by retrieving velocity fields from the reference stream function.

\begin{figure}[H]
    \begin{center}
        \includesvg[width=\textwidth]{assets/images/qg3l_all_rmses.svg}
    \end{center}
    \caption{nRMSE on three-layer QG surface stream function (top row) and velocities (bottom row), averaged over 12 20-days periods, each separated by 20 days, on Z1 (left) and Z2 (right). Lines represent the nRMSE average and the shaded area the standard deviation. Models are identified by the line color, black for RG, brown for RG-F and blue for RG-SI. The dashed line represents non-regularized RG-SI (hence $\gamma = 0$ in~\cref{def:RG-SI-reg}). Y-axis is shared in each row.}\label{fig:qg3L-nrmses}
\end{figure}

The RG-SI method is able to reconstruct the surface stream function with high fidelity (averaged nRMSE on Z1 is about 0.029 and about 0.016 on Z2 as detailed in~\cref{tab:qg3l-nrmses}) along with its gradient (providing the velocity). The introduction of the regularization, not only allowing the variational problem to be well constrained, also halves the average nRMSE over both zones. Looking at the velocity errors, especially in Z2, it also clearly improves the reconstruction of the stream function gradient. In all four panels we also see that the regularization flattens the nRMSE, suggesting that it helps to constrain the model when observations are sparse.

\begin{table}[H]
    \centering
        \begin{tabular}{clcccc}
            & & RG-SI & RG-SI ($\gamma = 0$) & RG-F & RG0 \\
        \hline
            \multirow{2}{*}{nRMSE} &\multicolumn{1}{l}{Z1} & 0.029 & 0.060 & 0.15 & 0.38 \\
            &\multicolumn{1}{l}{Z2} & 0.016 & 0.040 & 0.13 & 0.38 \\
        \addlinespace
            \multirow{2}{*}{$\mathrm{nRMSE}_{\mathrm{vel}}$} & \multicolumn{1}{l}{Z1} & 0.094 & 0.18 & 0.41 & 0.71 \\
            & \multicolumn{1}{l}{Z2} & 0.061 & 0.22 & 0.71 & 1.24 \\
        \hline
    \end{tabular}
\caption{Averaged nRMSE values.}\label{tab:qg3l-nrmses}
\end{table}

Comparison with other methods shows that RG-SI outperforms both RG0 and RG-F. Such result highlights that the introduction of an external forcing (the RG-F methodology), although able to mitigate some of the RG-truncation error, is not as effective as the truncation-aware method as introduced in RG-SI.
It is also worth mentioning that, as a comparison, we have tried fitting $\tilde{\sf}_2$ using the reference $\sf_2$ field of the QG3L model. Plugging the resulting parameters into the RG-SI model yields similar results as RG-SI (corresponding figure can be found in~\cref{fig:qg3L-proj-rmses}, averaged nRMSE values over stream function are of 0.041 in Z1 and 0.012 in Z2, and of 0.14 and 0.044 regarding velocity).

The visual comparison of vorticity fields shows almost identical fields. An example of such fields after 5 and 20 days of simulation are displayed in~\cref{fig:qg3L-vorticity}

\begin{figure}[H]
    \begin{center}
        \includesvg[width=\textwidth]{assets/images/qg3l_vorticity_all_32_96_256_384.svg}
    \end{center}
    \caption{Surface vorticity fields retrieved from the first assimilation cycle performed in Z1, after variational minimization has been completed. The first column displays reference vorticity field while second, third, fourth and fifth show RG-SI, non-regularized RG-SI, RG-F and RG0 reconstructions, respectively. Top row displays vorticity fields after 5 days and bottom one after 20 days.}\label{fig:qg3L-vorticity}
\end{figure}

A careful comparison of regularized (second column) and non-regularized (third column) vorticity fields obtained with the RG-SI framework shows that the regularization improves small scale reconstruction with the visible result of a smoothed out vorticity field.


\subsubsection{Estimation of the sublayer stream function}

As the RG-SI method provides an estimation for the sublayer flow, $\sf_2$, we have evaluated the reconstruction of $\sf_2$ by $\tilde{\sf}_2$. Given the reference sublayer stream function from the three-layer QG simulation, nRMSE is evaluated for $\tilde{\sf}_2$. Averages and standard deviation are computed over the variational assimilation periods as for the surface stream function.

\begin{figure}[H]
    \begin{center}
        \includesvg[width=\textwidth]{assets/images/qg3l_psi2_rmses.svg}
    \end{center}
    \caption{nRMSE on three-layer QG sublayer stream function (averaged over 12 20-days periods, each separated by 20 days) on Z1 (left) and Z2 (right). RG-SI model is identified by the solid blue line and non-regularized RG-SI by the dashed blue line. Y-axis is shared by both panels.}\label{fig:qg3L-psi2-rmse}
\end{figure}

Although results show that the regularization improves the reconstruction of $\sf_2$ by $\tilde{\sf}_2$, the quality of the overall reconstruction is far from being as good as it was for the surface flow. We understand this result as follows: although the parameterized field is indeed $\sf_2$, the actual field is never directly used by the model. What matters when emulating the dynamics is $\Dt*{\tilde{\sf}_2}$. Therefore, although $\sf_2$ might not be accurately represented by $\tilde{\sf}_2$, as long as $\Dt*{\sf_2}$ is by $\Dt*{\tilde{\sf}_2}$, the surface dynamics can still be precisely reconstructed. Evaluating the nRMSE of $\Dt*{\tilde{\sf}_2}$ yields better results, as it is mostly stable around 0.5 (resp. 0.15) in Z1 (resp. Z2). Corresponding figures can be found in the appendix (\cref{fig:qg3L-Dtpsi2-rmse}). Although still not as good as top layer reconstruction scores, $\Dt*{\tilde{\sf_2}}$ is better reconstructed than the field itself. Actually, when projecting the reference sublayer stream function onto the reduced-order space, the resulting error in $\Dt*{\tilde{\sf_2}}$ is not very far from the one obtained using RG-SI.

\subsubsection{Perturbations of stratification parameters}

To assess the relevancy of the deformation radius correction presented in~\cref{sec:deformation-radius-correction}, we keep the same parameters for the three-layer QG reference model, and we modify the other models' ones. Stratification perturbation is prescribed by modifying the layer reference thickness of the models. New parameters and the resulting deformation radius are detailed in the following table

\begin{table}[H]
    \centering
    \begin{tabular}{lcccc}
        & QG3L & RG-SI & RG-F & RG0 \\
        \hline
        $H_1$ & $\qty{400}{\meter}$ & $\qty{150}{\meter}$ & $\qty{135}{\meter}$ & $\qty{150}{\meter}$ \\
        $H_2$ & $\qty{1100}{\meter}$ & $\qty{1350}{\meter}$ &  &  \\
        $H_3$ & $\qty{2600}{\meter}$ &  &  &  \\
        \midrule
        $R_d$ & $\qty{41}{\kilo\meter}$ & $\qty{19}{\kilo\meter}$ & $\qty{19}{\kilo\meter}$ & $\qty{21}{\kilo\meter}$ \\
        \hline
    \end{tabular}
    \caption{Model parameters in the perturbed framework. Other model parameters were retained as shown in~\cref{tab:qg3l-model-parameters}.}\label{tab:qg3l-modified-model-parameters}
\end{table}

We then conducted the same experiment as before and monitored the ability of the different models to tamper incorrect parameters and reconstruct the surface dynamics.

\begin{figure}[H]
    \begin{center}
        \includesvg[width=\textwidth]{assets/images/qg3l_pert_rmses.svg}
    \end{center}
    \caption{nRMSE on three-layer QG surface stream function and velocities (average over 12 different assimilation periods each of 20 days) on Z1 (left) and Z2 (right) using perturbed models. Lines represent the nRMSE average and the shaded area the standard deviation. Models are identified by the line color, black for RG, brown for RG-F and dark blue for RG-SI. A light blue line is used for the RG-SI model with $\alpha=0$ (hence no deformation radius correction).}
\end{figure}

Here again, RG-SI accurately reconstructs the surface flow dynamics. The comparison of error scores of the standard RG-SI model and its counterpart with fixed $\alpha = 0$ shows that the introduction of a single additional parameter into the control vector results in an error reduction of more than 2 (see~\cref{tab:qg3l-pert-nrmses}). $\mathrm{nRMSE}_{\mathrm{vel}}$ scores are not displayed as they do not provide additional information to the nRMSE profiles.

\begin{table}[H]
    \centering
        \begin{tabular}{clcccc}
            & & RG-SI & RG-SI ($\alpha = 0$) & RG-F & RG0 \\
        \hline
            \multirow{2}{*}{nRMSE} & Z1 & 0.034 & 0.072 & 0.22 & 0.40 \\
                                   & Z2 & 0.011 & 0.027 & 0.14 & 0.41 \\
        \hline
    \end{tabular}
\caption{Averaged nRMSE values.}\label{tab:qg3l-pert-nrmses}
\end{table}

Relative comparison with the error scores from~\cref{tab:qg3l-nrmses} shows that, overall, RG-SI still performs worse in Z1 in the perturbed framework ($\sim 17 \% $ relative error degradation) and, surprisingly performs better in Z2 ($\sim 31\%$ improvement). RG-F performances, however, are both degraded in the perturbed framework, notably in Z1 ($\sim 47 \%$ relative degradation) and slightly in Z2 ($\sim 8 \%$ degradation). Nonetheless, perturbed RG-SI still outperform non-perturbed RG-F in reconstructing the surface flow of the three-layer QG model.

\subsection{Validation with realistic OSSE}

Using surface fields from the eNATL60 simulation~\cite{brodeau2020}, we monitor RG-SI's reconstruction of a fluid with ageostrophic dynamics. We have focused on a single area of dimension $\qty{680}{\kilo\meter}$ x $\qty{1030}{\kilo\meter}$ centered on $\qty{55}{\degree}$W $\qty{40}{\degree}$N.
Most parameters for all models are the exact same as the ones from~\cref{tab:qg3l-model-parameters}, apart from the time step, set to $\dt = \qty{3600}{\second}$ for all the models and the Coriolis parameters which were set to $f_0 = \qty{9.373}{\per\second}$ and $\beta = \qty{1.753}{\per\meter\per\second}$. Surface stream function anomaly is inferred from eNATL60's SSH and atmospheric pressures anomalies $p_{\mathrm{atm}}$, the latter being obtained from the Drakkar Forcing Set 5.2~\cite{dussin2016}. Anomalies are computed by removing its spatial average from the given field. Therefore, introducing the space-average operator ${\langle \cdot \rangle}_{XY}$, surface stream function anomaly $\Psi_{\mathrm{surf}}$ reads 

\begin{equation}\label{def:sf-processing}
    \Psi_{\mathrm{surf}} = \frac{p_{\mathrm{atm}}}{f_0 \rho_0} + \frac{g}{f_0} SSH - \left(\frac{{\langle p_{\mathrm{atm}} \rangle}_{XY}}{f_0 \rho_0} + \frac{g}{f_0}{\langle SSH \rangle}_{XY}\right),
\end{equation}

where $\rho_0$ is the reference density for water. Additional experiments were conducted by removing the time-space average of both SSH and atmospheric pressure fields and yielded overall similar results. However, in such configurations, the reference fields contained large-scale spatial fluctuations corresponding to the ocean response to atmospheric pressure. Since the QG models are not able to capture such fluctuations, we decided to mitigate them by removing the space-averaged (instead of the time-space average) stream function components, as illustrated by~\cref{def:sf-processing}.

As the resulting stream function contain small-scale features, it is then smoothed using a Gaussian filter and then downsampled from the high-resolution, curvilinear coordinates of the eNATL60 simulation into the regular, Cartesian grid of the QG models. It is this smoothed, down-scaled field that is used for both initial and boundary conditions of the QG models. The observations however are sampled from the non-filtered surface stream function field, using the same OSSE as for the three-layer QG comparison. A total of 10333 pixels are observed during a single assimilation cycle. Control vectors sizes are given in~\cref{tab:enatl60-control-vector-size}.

\begin{table}[H]
    \centering
    \begin{tabular}{lccc}
          & RG-SI & RG-F & RG0 \\
        \hline
        Control vector size & 5475 & 6608 & 0 \\
        \hline
    \end{tabular}
    \caption{Comparison of the control vector size for the different models in the eNATL60 experiment.}\label{tab:enatl60-control-vector-size}
\end{table}

Once again variational assimilation is conducted over 20-days periods and the quality of the reconstruction is assessed using stream function and velocity nRMSEs to compare reference fields to their reconstructed counterparts. Assimilation is conducted over 4 consecutive periods, during two different seasons, summer (July, August, September) and winter (January, February, March). Results are shown in~\cref{fig:enatl60-nrmse}. Velocity nRMSE is evaluated using surface zonal and meridional velocities from eNATL60 simulation as references.

\begin{figure}[H]
    \begin{center}
        \includesvg[width=\textwidth]{assets/images/enatl60_all_rmses.svg}
    \end{center}
    \caption{nRMSE on eNATL60 surface stream function (top row) and velocities (bottom row), averaged over 4 consecutive 20-days periods. Summer is displayed in the left column and winter in the right one. Models are identified by the line color, black for RG, brown for RG-F and blue for RG-SI. Solid line represents the averaged over the 4 periods and the shaded area the empirical standard deviation. Y-axis is shared in each row.}\label{fig:enatl60-nrmse}
\end{figure}

Despite the complexity of the realistic flow, RG-SI is still able to reconstruct surface dynamics with precision. The stream function nRMSE averages at 0.13 (resp. 0.11) for summer (resp. winter), demonstrating the ability of the model in reconstructing surface stream function. Regarding velocity, nRMSE scores remain significant (see~\cref{tab:enatl60-nrmses}). Indeed, initial velocity retrieved from initial condition already conveys a significant error (around 0.2-0.25), reflecting the limitation of our method in emulating ageostrophic components or small-scale patterns of the surface flow.

\begin{table}[H]
    \centering
        \begin{tabular}{clccc}
            & & RG-SI & RG-F & RG0 \\
        \hline
            \multirow{2}{*}{nRMSE} & Summer & 0.13 & 0.19 & 0.46 \\
            & Winter & 0.11 & 0.15 & 0.35 \\
        \addlinespace
            \multirow{2}{*}{$\mathrm{nRMSE}_{\mathrm{vel}}$} &  Summer & 0.41 & 0.51 & 0.87 \\
            &  Winter & 0.43 & 0.49 & 0.74  \\
        \hline
    \end{tabular}
\caption{Averaged nRMSE values.}\label{tab:enatl60-nrmses}
\end{table}

Such results indicate that a great component of the error conducted by the RG model can be explained by a field advected by geostrophic flow (whether only the sublayer geostrophic stream function or a more complex field). Regardless of the season (summer or winter), RG-SI provides an improvement over the RG-F method and significantly decreases the error that was conducted by the RG0 model.

The visual comparison of the surface vorticity field demonstrates the ability of the RG-SI method in reconstructing a realistic flow. An example of the vorticity fields is shown in~\cref{fig:enatl60-vorticity}. The overall comparison of the models show that RG-SI is able to reconstruct larger vortices, although its resolution prevents it from rendering smaller ones. RG-SI also seems to outperform the other models in locating the larger vortices.

\begin{figure}[H]
    \begin{center}
        \includesvg[width=\textwidth]{assets/images/enatl60_vorticity_summer.svg}
    \end{center}
    \caption{Surface vorticity fields retrieved from first assimilation cycle performed in summer, after variational minimization has been completed. First column displays reference vorticity field while second, third and fourth show RG-SI, RG-F and RG0 reconstructions, respectively. Top row displays vorticity field after 5 days and bottom one after 20 days.}\label{fig:enatl60-vorticity}
\end{figure}

\section{Conclusions and perspectives}

In order to improve sea surface reconstruction from altimetric observations using a reduced-gravity (RG) model, we have introduced a new formulation for the model error. This method emulates the subsurface stream function to account for the multilayer-truncation error that inherently comes with the RG assumption. Benefiting from such formulation, it is possible to express subsurface PV along with a diagnostic evaluation of its conservation.
We formulated a variational assimilation framework using such model, and we replaced the usual background regularization by a physics-based regularization enforcing subsurface PV conservation. Using OSSE-sampled surface data from a three-layer QG simulation, we reconstructed surface state through the minimization of the resulting cost function. Both nRMSE scores and visual comparison of the surface vorticity fields showed the benefits of using such error parameterization and regularization: stream function nRMSE showed a 80 to 85\% error reduction when compared to a model corrected by an external forcing. These findings confirm the ability of the method in mitigating truncation error thus highlighting the benefits of refining error parameterization using theoretical knowledge of the dynamics.
However, despite being able to provide an estimation for the subsurface stream function, comparison to actual subsurface stream function in the three-layer QG experiment yielded moderate results. These results can be explained by the fact that, when correcting surface advection equation, only the subsurface stream function's material derivative is accessed and is therefore better constrained by observations than the actual field is.
To evaluate our method over richer dynamics we then replicated the same experiment with observations sampled from the eNATL60 simulation. Results analysis has shown that the RG-SI methodology was able to improve surface reconstruction even for non-QG flows yet with a more modest improvement (25 to 30\% in stream function nRMSE). Nonetheless, although our method makes a strong assumption on the dynamics of the flow (as it enforces multilayer QG dynamics), such improvement over realistic flow suggest that a significant part of the reconstruction error can be explained by a field advected by the geostrophic flow. However, its identification to the subsurface stream function, in the context of a realistic flow, can be discussed.


Despite the improvement made over other reconstruction methods, the reconstruction error in the eNATL60 experiment remains significant (more than 0.1 for stream function nRMSE and above 0.4 for velocity error), suggesting that there might be some benefits in enriching the method. Indeed, dynamics is not purely QG anymore, and contains ageostrophic features (not accounted for in the QG framework).
Hence, rather than using the QG PV expression to formulate the error term, it might be interesting to dig into the next-order terms of the QG approximation~\cite{muraki1999}. Doing so would restore the proximity with Ertel's potential vorticity~\cite{muller1995} thus favoring the emergence of ageostrophic dynamics. 
Additionally, since the error parameterization, as formulated in RG-SI, only accounts for a specific source of error, the integration of physical forcing into the model might help better approximate the surface flow. Bottom topography for instance, cannot be taken into account by the RG model (as pointed out by~\citeA{hurlburt1986}) but could be integrated into the RG-SI methodology through appropriate modifications of the regularization term.
Hybridizing RG-SI and RG-F methods could allow for the representation of more generic error by confining RG-SI's error mitigation scope to truncation and leaving other unresolved process to be accounted for by the external forcing of RG-F term.

As all the presented results were established using simulated data, the next step would be to apply the RG-SI method to real observation. Although the eNATL60 simulation is realistic enough to grant confidence into the model's ability to handle real data, moving to such data yields notable challenges. Especially, initial and boundary conditions on the surface stream function would have to be addressed, whether estimated through climatology or gridded products, or included into the control vector. The incorporation of the dynamics of certain tracers (sea surface temperature for example) would also allow for using others sources of data. By doing so it would be possible to increase the coverage of the data and therefore to refine surface state reconstruction.

This new method highlights the benefits of deriving model error parameterization from the general physical framework of the model. The use of physics-based regularization has also proven to be relevant, especially in the absence of any prior estimate of the parameterized field, which can be very useful when there are no direct observations of the corresponding field. Successfully applying RG-SI to real SWOT data would be a significant step forward, assessing the relevancy of the method in a real-world application.


%%

%  Numbered lines in equations:
%  To add line numbers to lines in equations,
%  \begin{linenomath*}
%  \begin{equation}
%  \end{equation}
%  \end{linenomath*}



%% Enter Figures and Tables near as possible to where they are first mentioned:
%
% DO NOT USE \psfrag or \subfigure commands.
%
% Figure captions go below the figure.
% Acronyms used in figure captions will be spelled out in the final, published version.

% Table titles go above tables;  other caption information
%  should be placed in last line of the table, using
% \multicolumn2l{$^a$ This is a table note.}
% NOTE that there is no difference between table caption and table heading in the final, published version
%
%----------------
% EXAMPLE FIGURES
%
% \begin{figure}
% \includegraphics{example.png}
% \caption{caption}
% \end{figure}
%
% Giving latex a width will help it to scale the figure properly. A simple trick is to use \textwidth. Try this if large figures run off the side of the page.
% \begin{figure}
% \noindent\includegraphics[width=\textwidth]{anothersample.png}
%\caption{caption}
%\label{pngfiguresample}
%\end{figure}
%
%
% If you get an error about an unknown bounding box, try specifying the width and height of the figure with the natwidth and natheight options. This is common when trying to add a PDF figure without pdflatex.
% \begin{figure}
% \noindent\includegraphics[natwidth=800px,natheight=600px]{samplefigure.pdf}
%\caption{caption}
%\label{pdffiguresample}
%\end{figure}
%
%
% PDFLatex does not seem to be able to process EPS figures. You may want to try the epstopdf package.
%

%
% ---------------
% EXAMPLE TABLE
%
% \begin{table}
% \caption{Time of the Transition Between Phase 1 and Phase 2$^{a}$}
% \centering
% \begin{tabular}{l c}
% \hline
%  Run  & Time (min)  \\
% \hline
%   $l1$  & 260   \\
%   $l2$  & 300   \\
%   $l3$  & 340   \\
%   $h1$  & 270   \\
%   $h2$  & 250   \\
%   $h3$  & 380   \\
%   $r1$  & 370   \\
%   $r2$  & 390   \\
% \hline
% \multicolumn{2}{l}{$^{a}$Footnote text here.}
% \end{tabular}
% \end{table}

%%%%%%%%%%%%%%%%%%%%%%%%%%%%%%%%%%%%%%%%%%%%%%%
% SIDEWAYS FIGURES and TABLES
% AGU prefers the use of {sidewaystable} over {landscapetable} as it causes fewer problems.
%
% \begin{sidewaysfigure}
% \includegraphics[width=20pc]{figsamp}
% \caption{caption here}
% \label{newfig}
% \end{sidewaysfigure}
%
%  \begin{sidewaystable}
%  \caption{Caption here}
% \label{tab:signif_gap_clos}
%  \begin{tabular}{ccc}
% one&two&three\\
% four&five&six
%  \end{tabular}
%  \end{sidewaystable}

%% If using numbered lines, please surround equations with \begin{linenomath*}...\end{linenomath*}
%\begin{linenomath*}
%\begin{equation}
%y|{f} \sim g(m, \sigma),
%\end{equation}
%\end{linenomath*}

%%% End of body of article

%%%%%%%%%%%%%%%%%%%%%%%%%%%%%%%%%%%%%%%%%%%%%%%
%% Optional Appendices go here
%
% The \appendix command resets counters and redefines section heads
%
% After typing \appendix
%
%\section{Here Is Appendix Title}
% will show
% A: Here Is Appendix Title
%
%\appendix
%\section{Here is a sample appendix}

\newpage

\appendix

\section{RG-SI sublayer parameterization}\label{appendix:rg-si-parameterization}

Gaussian envelopes' standard deviations are chosen so that consecutive envelopes cross at $\sqrt{2}$. Hence, considering consecutive time-related (resp., space-related) envelopes to be distant from $\Delta t$ (resp., $\Delta x$ and $\Delta y$), their standard deviations are given by

\begin{subequations}
    \begin{align}
        \sigma_T & = \frac{\Delta t}{\sqrt{2\ln{2}}}, \\
        \sigma_X & = \frac{\Delta x}{\sqrt{2\ln{2}}}, \\
        \sigma_Y & = \frac{\Delta y}{\sqrt{2\ln{2}}}.
    \end{align}
\end{subequations}

\section{RG-F implementation}\label{appendix:rg-f-implementation}

RG-F's forcing is parameterized using a wavelet-like decomposition. Support functions span over different scales, both in time and space, and are built so that large spatial patterns are associated to large temporal patterns. Implementation is inspired from the methodology of~\cite{leguillou2025}. Wavelet functions are defined using cosine and sine functions, windowed using a Gaussian envelope (both in time and space). Since these cosine and sine functions are associated with different wave numbers, envelope widths are adjusted accordingly (standard deviation is computed from the inverse of the corresponding wave number). Lower frequency functions are thus windowed by broader envelopes.

Regarding the regularization, we penalize deviation from the 0 for the forcing. Supposing that we have considered $N$ different scales, since each decomposition coefficient weight a single basis function, it can be matched to its corresponding scale. Regrouping the coefficient by similar scale yields $N$ sets of coefficients, $\left( \theta_1 \in \mathbb{R}^{M_1}, \dots, \theta_N \in \mathbb{R}^{M_N} \right)$. We have tried to equally penalize these coefficients by defining

\begin{equation}
    \loss{reg}\left(\theta\right) = \sum_{i=1}^{N} {\lVert \theta_{i} \rVert}^2,
\end{equation}

but such penalization was actually preventing the model from accurately reconstructing the surface field. Indeed, forcing was either too strongly constrained and thus hardly able to improve the RG0 model, or too loosely constrained and producing small-scale-intensified fields that resulted in very noisy surface stream function reconstructions. To overcome this issue, we penalize the coefficients according to their scale. Hence, denoting $\left( k_1, \dots, k_N \right)$ the typical wave numbers of each scale, we penalize the decomposition coefficients using


\begin{equation}
    \loss{reg}\left(\theta\right) = \sum_{i=1}^{N} {\left( \frac{k_i}{k_1} \right)}^{2}{\lVert \theta_{i} \rVert}^2.
\end{equation}

Indeed, since small scales (large $k_i$) are expected to carry less energy, we increase the strength of the regularization for such scales following a 5/3 power law. Using this regularization has improved the RG-F reconstruction capacities while producing smoother fields. Relative weighting of the regularization cost function in comparison to the observation loss is then adjusted using the hyperparameter $\gamma$ as introduced in~\cref{def:RG-SI-reg}.


\section{Projected three-layer QG sublayer stream function performances}

\begin{figure}[H]
    \begin{center}
        \includesvg[width=\textwidth]{assets/images/qg3l_all_rmses_.svg}
    \end{center}
    \caption{nRMSE on surface stream function (top row) and velocities (bottom row),averaged over 12 20-days periods, each separated by 20 days, on Z1 (left) and Z2 (right). Standard RG-SI model is identified by the solid blue line, RG-SI with coefficient obtained through the projection of $\sf_2$ (from the reference three-layer QG model) onto the reduced order space is shown by the green line and RG-F by the brown one. Y-axis is shared by both panels.}\label{fig:qg3L-proj-rmses}
\end{figure}

\begin{figure}[H]
    \begin{center}
        \includesvg[width=\textwidth]{assets/images/qg3l_Dt_psi2_rmses.svg}
    \end{center}
    \caption{nRMSE on material derivative of three-layer QG sublayer stream function (averaged over 12 20-days periods, each separated by 20 days) on Z1 (left) and Z2 (right). Standard RG-SI model is identified by the solid blue line whereas RG-SI with coefficient obtained through the projection of $\sf_2$ (from the reference three-layer QG model) onto the reduced order space is shown by the green line. Y-axis is shared by both panels.}\label{fig:qg3L-Dtpsi2-rmse}
\end{figure}

%%%%%%%%%%%%%%%%%%%%%%%%%%%%%%%%%%%%%%%%%%%%%%%
% Optional Glossary, Notation or Acronym section goes here:
%
% Glossary is only allowed in Reviews of Geophysics
%  \begin{glossary}
%  \term{Term}
%   Term Definition here
%  \term{Term}
%   Term Definition here
%  \term{Term}
%   Term Definition here
%  \end{glossary}


%%%%%%%%%%%%%%%%%%%%%%%%%%%%%%%%%%%%%%%%%%%%%%%
% Acronyms
%% NOTE that acronyms in the final published version will be spelled out when used in figure captions.
%   \begin{acronyms}
%   \acro{Acronym}
%   Definition here
%   \acro{EMOS}
%   Ensemble model output statistics
%   \acro{ECMWF}
%   Centre for Medium-Range Weather Forecasts
%   \end{acronyms}


%%%%%%%%%%%%%%%%%%%%%%%%%%%%%%%%%%%%%%%%%%%%%%%
% Notation
%   \begin{notation}
%   \notation{$a+b$} Notation Definition here
%   \notation{$e=mc^2$}
%   Equation in German-born physicist Albert Einstein's theory of special
%  relativity that showed that the increased relativistic mass ($m$) of a
%  body comes from the energy of motion of the body—that is, its kinetic
%  energy ($E$)—divided by the speed of light squared ($c^2$).
%   \end{notation}




%%%%%%%%%%%%%%%%%%%%%%%%%%%%%%%%%%%%%%%%%%%%%%%
%
% DATA SECTION and ACKNOWLEDGMENTS
%
%%%%%%%%%%%%%%%%%%%%%%%%%%%%%%%%%%%%%%%%%%%%%%%

\section*{Open Research Section}

Python codes used to emulate the models are available at \url{https://github.com/grigaut/QG/}.

\section*{Conflict of Interest declaration}

The authors declare there are no conflicts of interest for this manuscript.


%%%%%%%%%%%%%%%%%%%%%%%%%%%%%%%%%%%%%%%%%%%%%%%
% REFERENCES and BIBLIOGRAPHY
%
% \bibliography{<name of your .bib file>} don't specify the file extension
% don't specify bibliographystyle
%
%%%%%%%%%%%%%%%%%%%%%%%%%%%%%%%%%%%%%%%%%%%%%%%

\bibliography{refs}



%Reference citation instructions and examples:
%
% Please use ONLY \cite and \citeA for reference citations.
% \cite for parenthetical references
% ...as shown in recent studies (Simpson et al., 2019)
% \citeA for in-text citations
% ...Simpson et al. (2019) have shown...
%
%
%...as shown by \citeA{jskilby}.
%...as shown by \citeA{lewin76}, \citeA{carson86}, \citeA{bartoldy02}, and \citeA{rinaldi03}.
%...has been shown \citeA{jskilbye}.
%...has been shown \citeA{lewin76,carson86,bartoldy02,rinaldi03}.
%... \cite <i.e.>[]{lewin76,carson86,bartoldy02,rinaldi03}.
%...has been shown by \cite <e.g.,>[and others]{lewin76}.
%
% apacite uses < > for prenotes and [ ] for postnotes
% DO NOT use other cite commands (e.g., \citet, \citep, \citeyear, \nocite, \citealp, etc.).
%



\end{document}


More Information and Advice:

%%%%%%%%%%%%%%%%%%%%%%%%%%%%%%%%%%%%%%%%%%%%%%%
%
%  SECTION HEADS
%
%%%%%%%%%%%%%%%%%%%%%%%%%%%%%%%%%%%%%%%%%%%%%%%

% Capitalize the first letter of each word (except for
% prepositions, conjunctions, and articles that are
% three or fewer letters).

% AGU follows standard outline style; therefore, there cannot be a section 1 without
% a section 2, or a section 2.3.1 without a section 2.3.2.
% Please make sure your section numbers are balanced.
% ---------------
% Level 1 head
%
% Use the \section{} command to identify level 1 heads;
% type the appropriate head wording between the curly
% brackets, as shown below.
%
%An example:
%\section{Level 1 Head: Introduction}
%
% ---------------
% Level 2 head
%
% Use the \subsection{} command to identify level 2 heads.
%An example:
%\subsection{Level 2 Head}
%
% ---------------
% Level 3 head
%
% Use the \subsubsection{} command to identify level 3 heads
%An example:
%\subsubsection{Level 3 Head}
%
%---------------
% Level 4 head
%
% Use the \subsubsubsection{} command to identify level 3 heads
% An example:
%\subsubsubsection{Level 4 Head} An example.
%
%%%%%%%%%%%%%%%%%%%%%%%%%%%%%%%%%%%%%%%%%%%%%%%
%
%  IN-TEXT LISTS
%
%%%%%%%%%%%%%%%%%%%%%%%%%%%%%%%%%%%%%%%%%%%%%%%
%
% Do not use bulleted lists; enumerated lists are okay.
% \begin{enumerate}
% \item
% \item
% \item
% \end{enumerate}
%
%%%%%%%%%%%%%%%%%%%%%%%%%%%%%%%%%%%%%%%%%%%%%%%
%
%  EQUATIONS
%
%%%%%%%%%%%%%%%%%%%%%%%%%%%%%%%%%%%%%%%%%%%%%%%

% Single-line equations are centered.
% Equation arrays will appear left-aligned.

Math coded inside display math mode \[ ...\]
 will not be numbered, e.g.,:
 \[ x^2=y^2 + z^2\]

 Math coded inside \begin{equation} and \end{equation} will
 be automatically numbered, e.g.,:
 \begin{equation}
 x^2=y^2 + z^2
 \end{equation}


% To create multiline equations, use the
% \begin{eqnarray} and \end{eqnarray} environment
% as demonstrated below.
\begin{eqnarray}
  x_{1} & = & (x - x_{0}) \cos \Theta \nonumber \\
        && + (y - y_{0}) \sin \Theta  \nonumber \\
  y_{1} & = & -(x - x_{0}) \sin \Theta \nonumber \\
        && + (y - y_{0}) \cos \Theta.
\end{eqnarray}

%If you don't want an equation number, use the star form:
%\begin{eqnarray*}...\end{eqnarray*}

% Break each line at a sign of operation
% (+, -, etc.) if possible, with the sign of operation
% on the new line.

% Indent second and subsequent lines to align with
% the first character following the equal sign on the
% first line.

% Use an \hspace{} command to insert horizontal space
% into your equation if necessary. Place an appropriate
% unit of measure between the curly braces, e.g.
% \hspace{1in}; you may have to experiment to achieve
% the correct amount of space.


%%%%%%%%%%%%%%%%%%%%%%%%%%%%%%%%%%%%%%%%%%%%%%%
%
%  EQUATION NUMBERING: COUNTER
%
%%%%%%%%%%%%%%%%%%%%%%%%%%%%%%%%%%%%%%%%%%%%%%%

% You may change equation numbering by resetting
% the equation counter or by explicitly numbering
% an equation.

% To explicitly number an equation, type \eqnum{}
% (with the desired number between the brackets)
% after the \begin{equation} or \begin{eqnarray}
% command.  The \eqnum{} command will affect only
% the equation it appears with; LaTeX will number
% any equations appearing later in the manuscript
% according to the equation counter.
%

% If you have a multiline equation that needs only
% one equation number, use a \nonumber command in
% front of the double backslashes (\\) as shown in
% the multiline equation above.

% If you are using line numbers, remember to surround
% equations with \begin{linenomath*}...\end{linenomath*}

%  To add line numbers to lines in equations:
%  \begin{linenomath*}
%  \begin{equation}
%  \end{equation}
%  \end{linenomath*}



